

% --- Use Case U1.2 (U12 trong file) ---
\vspace{0.5cm}
\noindent
\begin{tabularx}{\textwidth}{|>{\raggedright\arraybackslash}p{4cm}|X|}
\hline
\textbf{Use-case ID} \newline \textbf{Use-case name} & U1.2 \newline \textbf{Xác nhận và ghi nhận ra vào tự động} \\ \hline
\textbf{Use-case overview} & Tự động ghi nhận lượt ra/vào bằng cách quét thẻ RFID và camera nhận diện mà không cần nhân viên thao tác.\\ \hline
\textbf{Actors} & 1. Khách gửi xe có tài khoản \newline 2. HCMUT\_DATACORE  \\ \hline
\textbf{Preconditions} & 
1. Thẻ sinh viên / giảng viên còn hiệu lực.\newline
2. Biển số xe đã được đăng ký trên hệ thống.  \\ \hline
\textbf{Trigger} & Đầu đọc thẻ tại trạm kiểm soát bắt được tín hiệu từ thẻ của khách. \\ \hline
\textbf{Steps} & 
1. Khách sử dụng thẻ sinh viên/giảng viên để quét tại lối vào/ra.\newline
2. Hệ thống đọc mã số sinh viên/giảng viên từ thẻ.\newline
3. Hệ thống đối chiếu tài khoản đã có trên HCMUT\_DATACORE chưa.\newline
4. Hệ thống kích hoạt camera chụp biển số xe.\newline
5. Hệ thống kiểm tra biển số có khớp với biển số đã được đăng ký.\newline
6. Nếu biển số hợp lệ, hệ thống ghi nhận lượt xe tương ứng và cập nhật vị trí trên Parking Map sau 5s cảm biến ổn định.\newline
7. Đối với khách không đăng ký gói, hệ thống tự động cộng dồn phí; đối với khách có gói còn thời hạn thì bỏ qua bước tính phí.  \\ \hline
\textbf{Post conditions} & Hệ thống ghi nhận lượt vào ra.  \\ \hline
\textbf{Exception flow} & 
- \textbf{E1 (Thẻ không hợp lệ):} Thẻ hết hạn hoặc bị khóa. Hệ thống báo lỗi trên màn hình LED, rào chắn không mở.\newline
- \textbf{E2 (Sai biển số):} Biển số không khớp với đăng ký. Hệ thống báo lỗi trên màn hình LED, rào chắn không mở.  \\ \hline
\end{tabularx}
\newpage

% --- Use Case U1.3 (U13 trong file) ---
\vspace{0.5cm}
\noindent
\begin{tabularx}{\textwidth}{|>{\raggedright\arraybackslash}p{4cm}|X|}
\hline
\textbf{Use-case ID} \newline \textbf{Use-case name} & U1.3 \newline \textbf{Xác nhận ra vào thủ công} \\ \hline
\textbf{Use-case overview} & Thanh toán phí đỗ xe trực tiếp bằng tiền mặt cho nhân viên tại cổng ra. \\ \hline
\textbf{Actors} & 1. Khách không có tài khoản \newline 2. Nhân viên vận hành bãi xe  \\ \hline
\textbf{Preconditions} & Khách có thẻ từ đã được cấp từ nhân viên bãi xe.  \\ \hline
\textbf{Trigger} & Khách điều khiển xe đến lối vào/ra.  \\ \hline
\textbf{Steps} & 
1. Nhân viên cấp thẻ từ cho khách khi vào bãi.\newline
2. Khách vào/ra bãi giữ xe.\newline
3. Hệ thống ghi nhận lượt và cập nhật vị trí xe trên Parking Map.\newline
4. Khi khách ra, nhân viên nhận thẻ từ và tiền mặt từ khách.\newline
5. Nhân viên quẹt thẻ để xác nhận. \\ \hline
\textbf{Post conditions} & Nhân viên xác nhận đã nhận tiền; hệ thống ghi nhận lượt vào ra.  \\ \hline
\textbf{Exception flow} & Không có.  \\ \hline
\end{tabularx}
\newpage

% --- Use Case U1.4-15 (U14-15 trong file) ---
\vspace{0.5cm}
\noindent
\begin{tabularx}{\textwidth}{|>{\raggedright\arraybackslash}p{4cm}|X|}
\hline
\textbf{Use-case ID} \newline \textbf{Use-case name} & U1.4-15 \newline \textbf{Đăng ký biển số xe} \\ \hline
\textbf{Use-case overview} & Khách hàng đăng ký biển số xe mới lên hệ thống và chờ nhân viên xác nhận để đảm bảo tính hợp lệ.  \\ \hline
\textbf{Actors} & 1. Khách hàng có tài khoản \newline 2. Nhân viên vận hành bãi xe  \\ \hline
\textbf{Preconditions} & Khách hàng đã đăng nhập thành công trên hệ thống.  \\ \hline
\textbf{Trigger} & Bấm nút "Đăng ký Biển số xe" trên giao diện web. \\ \hline
\textbf{Steps} & 
1. Khách hàng tải ảnh giấy tờ (Cà vẹt) định dạng .JPG và nhập biển số.\newline
2. Hệ thống lưu dữ liệu "Chờ xác nhận", đẩy thông báo về màn hình Nhân viên.\newline
3. Nhân viên đối chiếu ảnh với thông tin.\newline
4. Nhân viên nhấn "Xác nhận duyệt".\newline
5. Hệ thống cập nhật trạng thái "Đã duyệt" và thông báo cho Khách.  \\ \hline
\textbf{Post conditions} & Biển số xe được lưu, liên kết với tài khoản và chuyển sang trạng thái "Đã duyệt". \\ \hline
\textbf{Exception flow} & 
- \textbf{E1 (Lỗi định dạng):} Nhập thiếu, sai định dạng hoặc không khớp cà vẹt; hệ thống yêu cầu nhập lại.\newline
- \textbf{E2 (Biển số đã tồn tại):} Hệ thống chặn thao tác và báo lỗi.\newline
- \textbf{E3 (Từ chối duyệt):} Ảnh mờ hoặc giả mạo; nhân viên từ chối và nhập lý do.  \\ \hline
\end{tabularx}
\newpage

% --- Use Case U1.6 (U16 trong file) ---
\vspace{0.5cm}
\noindent
\begin{tabularx}{\textwidth}{|>{\raggedright\arraybackslash}p{4cm}|X|}
\hline
\textbf{Use-case ID} \newline \textbf{Use-case name} & U1.6 \newline \textbf{Đăng ký gói theo tháng} \\ \hline
\textbf{Use-case overview} & Khách hàng mua hoặc gia hạn gói cước gửi xe theo tháng qua cổng thanh toán tích hợp (BKPay). \\ \hline
\textbf{Actors} & 1. Khách hàng có tài khoản \newline 2. BKPay \newline 3. HCMUT\_SSO  \\ \hline
\textbf{Preconditions} & Khách đã đăng nhập thành công trên HCMUT\_SSO. \\ \hline
\textbf{Trigger} & Bấm nút "Đăng ký gói theo tháng" trên giao diện web.\\ \hline
\textbf{Steps} & 
1. Hệ thống tạo mã đơn, tính tổng tiền và chuyển hướng qua cổng BKPay.\newline
2. Khách thực hiện thanh toán trên giao diện BKPay.\newline
3. BKPay xử lý giao dịch và gửi tín hiệu thành công về hệ thống.\newline
4. Hệ thống kích hoạt gói cước và tính ngày hết hạn mới.\newline
5. Hệ thống hiển thị biên lai và thông báo thành công.  \\ \hline
\textbf{Post conditions} & Giao dịch thành công, gói vé tháng được kích hoạt/gia hạn. \\ \hline
\textbf{Exception flow} & 
- \textbf{E1 (Quá thời gian giao dịch):} Quá 5 phút tại trang BKPay; hệ thống tự động hủy đơn.  \\ \hline
\end{tabularx}

